\subsection{对偶向量空间}

设 \(\mathbf{F}\) 和 \(\mathbf{G}\) 是两个实向量空间，\((x, y) \mapsto \mathbf{B}(x, y)\) 是 \(\mathbf{F} \times \mathbf{G}\) 上的一个双线性型。我们说双线性型 \(\mathbf{B}\) 使向量空间 \(\mathbf{F}\) 和 \(\mathbf{G}\) 成对偶，或说 F 和 G 成对偶（关于 B）。回忆我们称 \(x \in F\) 和 \(y \in G\) （对由 B 定义的对偶）正交，如果 \(\mathbf{B}(x, y) = 0\) ；我们称 F 的一个子集 M 和 G 的一个子集 N 正交，如果每个 \(x \in M\) 与每个 \(y \in N\) 都正交 (A, IX, § 1.2)。

我们说由 B 定义的对偶在 F 中（相应地在 G 中）分离，如果它满足下列条件：

(D₁) 对 F 中每个 \(x \neq 0\) ，存在 \(y \in G\) 使 \(\mathrm{B}(x, y) \neq 0\) 。（相应地 (D \(_{II}\) ) 对 G 中每个 y ≠ 0，存在 x ∈ F 使 B(x, y) ≠ 0。）

由 B 定义的对偶称为分离的，如果它在 F 中和 G 中都分离。为此，必须且只需双线性型 B 在 A, IX, § 1.1 的意义下是分离的。更精确地，我们有下面的结果：

命题 1. --- 设 F, G 是两个实向量空间，B 是 F × G 上的一个双线性型。设

\[
d _ {\mathrm{B}}: y \mapsto \mathrm{B} (., y),
\]

\[
s _ {\mathbf {B}} \colon x \mapsto \mathrm{B} (x,.)
\]

分别是 G 到 F 的对偶 F* 中和 F 到 G 的对偶 G* 中的线性映射，它们分别伴随于 B 的右与左 (A, IX, § 1.1)。则 B 使 F 和 G 成一个在 G 中（相应地在 F 中）分离的对偶，必须且只需 \(d_{\mathbf{B}}\) （相应地 \(s_{\mathbf{B}}\) ）是单射。

当 F 和 G 被 B 置于分离对偶时，我们常把 F（相应地 G）等同于 \(G^{*}\) （相应地 \(F^{*}\) ）的一个子空间（借助 \(s_{B}\) （相应地 \(d_{B}\) ））。当我们把 F（相应地 G）看作 \(G^{*}\) （相应地 \(F^{*}\) ）的一个子空间而不指明这个等同如何实现时，我们总是在使用前述的等同；这时双线性型 B 等同于典范双线性型在 \(F \times G\) 上的限制：

\[
(x ^ {*}, x) \mapsto \langle x, x ^ {*} \rangle \quad (\text { resp. } (x, x ^ {*}) \mapsto \langle x, x ^ {*} \rangle).
\]

例. --- 1) 设 E 是一个向量空间，E* 是它的对偶。典范双线性型 \((x, x^{*}) \mapsto \langle x, x^{*} \rangle\) （\(E \times E^{*}\) 上的）(A, II, §2.3) 使 E 和 \(E^{*}\) 成分离对偶：因为 \((\mathrm{D}_{\mathrm{II}})\) 由关系 \(x^{*} \neq 0\) 的定义得出，另一方面我们知道，对 E 中所有 \(x \neq 0\) ，存在线性型 \(x^{*} \in E^{*}\) 使 \(\langle x, x^{*} \rangle \neq 0\) (A, II, §7.5, th. 6)，这就证明了 \((\mathrm{D}_{\mathrm{I}})\) ；这里 E 与 \(E^{**}\) 的一个子空间的等同由典范映射 \(c_{E}\) 实现 (loc. cit.)。

当 E 是有限维时，与 E 成分离对偶的 \(E^{*}\) 的子空间 G（借助典范双线性型在 \(E \times G\) 上的限制）只有 \(E^{*}\) 自身；因为这时 E 典范地等同于 \(E^{**}\) (loc. cit.)，倘若有 \(G \neq E^{*}\) ，就将存在 E 中的 \(a \neq 0\) ，使 \(\langle a, x^{*} \rangle = 0\) （对所有 \(x^{*} \in G\) ） (A, II, §7.5, th. 7)，这与假设矛盾。

\begin{enumerate}
\def\labelenumi{\arabic{enumi})}
\setcounter{enumi}{1}
\tightlist
\item
  当 E 是无限维向量空间且 \(E'\) 是 \(E^{*}\) 的一个向量子空间时，E 与 \(E'\) 之间的对偶（由典范双线性型在 \(E \times E'\) 上的限制所定义）总是在 \(E'\) 中分离的；即使 \(E' \neq E^{*}\) ，它也可以在 E 中分离。最重要的情形是 E 为拓扑向量空间的情形。
\end{enumerate}

定义 1. --- 拓扑向量空间 E 的对偶指的是向量空间 E 的对偶 E* 中由 E 上的连续线性型组成的子空间 E'。

当 E 是 Hausdorff 局部凸空间时，E 与其对偶 \(E'\) 之间的对偶是分离的：这由 Hahn-Banach 定理 (II, 第 24 页, cor. 1) 得出，即对 E 中每个 \(x \neq 0\) ，存在 \(x' \in E'\) 使 \(\langle x, x' \rangle \neq 0\) 。

注记. --- 1) 当 E 是拓扑向量空间时，向量空间 E 的对偶 E* 将称为 E 的代数对偶，以避免混淆。我们还注意，E* 是给 E 赋予最细局部凸拓扑 (II, 第 25 页, 例 2) 后所得拓扑向量空间的对偶。

\begin{enumerate}
\def\labelenumi{\arabic{enumi})}
\setcounter{enumi}{1}
\item
  拓扑向量空间的对偶 \(\mathbf{E}'\) 本身不带有拓扑，除非明确声明。
\item
  若 \(\mathbf{F}\) 与 \(\mathbf{G} \subset \mathbf{F}^*\) 由典范双线性型成分离对偶，则 \(\mathbf{F}\) 与 \(\mathbf{G}_1\) 也成分离对偶，对子空间 \(\mathbf{G}_1\) （\(\mathbf{F}^*\) 的，满足 \(\mathbf{G} \subset \mathbf{G}_1\) ）。
\end{enumerate}

